\section{Second-Degree Equations and Conic Sections}

Lines arise from equations in which the coordinates occur only to the first
power. The next natural class allows terms of total degree two.

\begin{definitionbox}
A \emph{second-degree equation} in the Cartesian plane has the form
\[
\boxed{Ax^2+Bxy+Cy^2+Dx+Ey+F=0},
\]
where $A,B,C,D,E,F\in\mathbb R$ and at least one of $A,B,C$ is nonzero. Its
solution set consists of all points $(x,y)$ satisfying the equation.
\end{definitionbox}

\begin{interpretationbox}[The role of the quadratic part]
The terms
\[
Ax^2+Bxy+Cy^2
\]
form the quadratic part. They largely determine the geometric type, while the
linear and constant terms mainly control position and scale.
\end{interpretationbox}

\subsection{A First Classification}

\begin{definitionbox}
For
\[
Ax^2+Bxy+Cy^2+Dx+Ey+F=0,
\]
define the conic discriminant
\[
\boxed{\Delta=B^2-4AC}.
\]
\end{definitionbox}

\begin{theorembox}
For a nondegenerate real conic,
\[
\begin{array}{c|c}
\Delta<0 & \text{ellipse type, including circles},\\
\Delta=0 & \text{parabola type},\\
\Delta>0 & \text{hyperbola type}.
\end{array}
\]
The sign is preserved by translations and rotations of the coordinate system.
\end{theorembox}

\begin{interpretationbox}[Why translations and rotations are used]
A translation moves the centre or vertex to a simpler position. A rotation
aligns the axes with the natural symmetry directions and can remove the mixed
term $Bxy$. These transformations simplify the description without changing the
underlying geometric curve.
\end{interpretationbox}

\begin{remarkbox}
The sign of $\Delta$ identifies the conic type only after degenerate cases are
separated. A second-degree equation may describe a point, no real points, one
line, or two lines.
\end{remarkbox}

\subsection{Circle}

\begin{definitionbox}
A circle with centre $(h,k)$ and radius $r>0$ has standard equation
\[
\boxed{(x-h)^2+(y-k)^2=r^2}.
\]
\end{definitionbox}

\begin{derivationbox}[Classification by the discriminant]
Expanding gives
\[
x^2+y^2-2hx-2ky+h^2+k^2-r^2=0.
\]
Here $A=C=1$ and $B=0$, so
\[
\Delta=-4<0.
\]
\end{derivationbox}

\begin{interpretationbox}[Circle as a special ellipse]
A circle is an ellipse-type conic whose two quadratic coefficients are equal and
whose natural axes have the same scale.
\end{interpretationbox}

\begin{center}
\begin{tikzpicture}[scale=0.82]
    \draw[axis] (-3.8,0) -- (4.2,0) node[right] {$x$};
    \draw[axis] (0,-3.3) -- (0,3.7) node[above] {$y$};
    \coordinate (C) at (1,0.6);
    \draw[subset] (C) circle (2);
    \fill[geometricSubsetColor] (C) circle (2pt) node[below right] {$(h,k)$};
    \draw[helper] (C) -- ++(35:2) node[midway,above] {$r$};
\end{tikzpicture}
\end{center}

\subsection{Ellipse}

\begin{definitionbox}
A standard ellipse centred at the origin has equation
\[
\boxed{
\frac{x^2}{a^2}+\frac{y^2}{b^2}=1
},
\qquad a>0,\quad b>0.
\]
The parameters $a$ and $b$ are the semi-axis lengths.
\end{definitionbox}

\begin{derivationbox}[Classification by the discriminant]
Multiplying by $a^2b^2$ gives
\[
b^2x^2+a^2y^2-a^2b^2=0.
\]
Thus $A=b^2$, $B=0$, and $C=a^2$, so
\[
\Delta=-4a^2b^2<0.
\]
\end{derivationbox}

\begin{remarkbox}
When $a=b$, the ellipse becomes a circle.
\end{remarkbox}

\begin{center}
\begin{tikzpicture}[scale=0.85]
    \draw[axis] (-4.2,0) -- (4.5,0) node[right] {$x$};
    \draw[axis] (0,-3.0) -- (0,3.3) node[above] {$y$};
    \draw[subset] (0,0) ellipse (3.1 and 1.8);
    \draw[helper] (0,0) -- (3.1,0) node[midway,below] {$a$};
    \draw[helper] (0,0) -- (0,1.8) node[midway,left] {$b$};
\end{tikzpicture}
\end{center}

\subsection{Parabola}

\begin{definitionbox}
A standard vertical parabola has equation
\[
\boxed{y=ax^2},
\qquad a\neq0.
\]
A translated form is
\[
\boxed{y-k=a(x-h)^2},
\]
with vertex $(h,k)$.
\end{definitionbox}

\begin{derivationbox}[Classification by the discriminant]
The equation $y=ax^2$ is equivalent to
\[
ax^2-y=0.
\]
Here $A=a$, $B=0$, and $C=0$, so
\[
\Delta=0.
\]
\end{derivationbox}

\begin{interpretationbox}[What the coefficient $a$ controls]
The sign of $a$ determines whether the parabola opens upward or downward, while
its magnitude controls how rapidly it widens.
\end{interpretationbox}

\begin{center}
\begin{tikzpicture}[scale=0.82]
    \draw[axis] (-3.8,0) -- (3.9,0) node[right] {$x$};
    \draw[axis] (0,-1.0) -- (0,5.2) node[above] {$y$};
    \draw[subset,domain=-2.2:2.2,samples=100] plot (\x,{0.85*\x*\x});
    \fill[geometricSubsetColor] (0,0) circle (2pt) node[below right] {vertex};
\end{tikzpicture}
\end{center}

\subsection{Hyperbola}

\begin{definitionbox}
A standard hyperbola centred at the origin has equation
\[
\boxed{
\frac{x^2}{a^2}-\frac{y^2}{b^2}=1
},
\qquad a>0,\quad b>0.
\]
Its asymptotes are
\[
\boxed{y=\frac ba x},
\qquad
\boxed{y=-\frac ba x}.
\]
\end{definitionbox}

\begin{derivationbox}[Classification by the discriminant]
Multiplying by $a^2b^2$ gives
\[
b^2x^2-a^2y^2-a^2b^2=0.
\]
Thus $A=b^2$, $B=0$, and $C=-a^2$, so
\[
\Delta=4a^2b^2>0.
\]
\end{derivationbox}

\begin{center}
\begin{tikzpicture}[scale=0.78]
    \draw[axis] (-5.0,0) -- (5.2,0) node[right] {$x$};
    \draw[axis] (0,-4.0) -- (0,4.2) node[above] {$y$};
    \draw[helper] (-4.6,-3.07) -- (4.6,3.07);
    \draw[helper] (-4.6,3.07) -- (4.6,-3.07);
    \draw[subset,domain=1.55:4.5,samples=100] plot (\x,{1.1*sqrt(\x*\x/2.25-1)});
    \draw[subset,domain=1.55:4.5,samples=100] plot (\x,{-1.1*sqrt(\x*\x/2.25-1)});
    \draw[subset,domain=-4.5:-1.55,samples=100] plot (\x,{1.1*sqrt(\x*\x/2.25-1)});
    \draw[subset,domain=-4.5:-1.55,samples=100] plot (\x,{-1.1*sqrt(\x*\x/2.25-1)});
\end{tikzpicture}
\end{center}

\subsection{The Mixed Term and Rotation}

\begin{corebox}[What $Bxy$ signals]
A mixed term indicates that the chosen Cartesian axes may not coincide with the
natural symmetry directions of the conic.
\end{corebox}

\begin{examplebox}
The equation
\[
x^2+2xy+y^2=1
\]
can be written as
\[
(x+y)^2=1.
\]
It is degenerate and describes the two parallel lines
\[
x+y=1,
\qquad
x+y=-1.
\]
\end{examplebox}

\begin{derivationbox}[Rotating coordinates]
Introduce new coordinates $(u,v)$ by
\[
x=u\cos\theta-v\sin\theta,
\qquad
y=u\sin\theta+v\cos\theta.
\]
Choosing $\theta$ appropriately removes the $uv$ term. When $A\neq C$, one may
choose
\[
\boxed{
\tan(2\theta)=\frac{B}{A-C}
}.
\]
The resulting equation can then be classified in the rotated coordinates using
standard forms.
\end{derivationbox}

\subsection{Degenerate Second-Degree Equations}

\begin{corebox}[A quadratic equation need not be a conic curve]
Factorisation or completing squares may reveal a point, no real points, one
line, or two lines instead of a nondegenerate circle, ellipse, parabola, or
hyperbola.
\end{corebox}

\begin{examplebox}
The equation
\[
x^2-y^2=0
\]
factors as
\[
(x-y)(x+y)=0.
\]
Its solution set is the union of two intersecting lines:
\[
y=x
\qquad\text{or}\qquad
y=-x.
\]
\end{examplebox}

\begin{center}
\begin{tikzpicture}[scale=0.75]
    \draw[axis] (-4.0,0) -- (4.2,0) node[right] {$x$};
    \draw[axis] (0,-3.8) -- (0,4.0) node[above] {$y$};
    \draw[subset] (-3.5,-3.5) -- (3.5,3.5);
    \draw[subset] (-3.5,3.5) -- (3.5,-3.5);
\end{tikzpicture}
\end{center}

\begin{examplebox}
Other common degenerate cases are
\[
x^2+y^2=0,
\]
which gives only $(0,0)$,
\[
x^2+y^2+1=0,
\]
which has no real solutions, and
\[
x^2=0,
\]
which describes the line $x=0$ counted algebraically twice.
\end{examplebox}

\begin{summarybox}[Recognition procedure]
Second-degree equations have the form
\[
Ax^2+Bxy+Cy^2+Dx+Ey+F=0.
\]
For nondegenerate real conics, $\Delta=B^2-4AC$ separates ellipse, parabola, and
hyperbola types. Translations and rotations reduce equations to standard forms,
while factorisation and completing squares expose degenerate cases.
\end{summarybox}
